\documentclass[border=4pt]{standalone}
\usepackage{graficas}

\begin{document}
\begin{tikzpicture}
    \begin{axis}[grafica, width=8cm, height=5cm, xmin=-1.8, xmax=2.6, ymin=-0.8, ymax=5.2]
        \addplot[red!80, dashed, line width=0.9pt, domain=-1.8:2.6] {x^2};
        % f(x) = (x^5-2x^2+3)/(x^3+1); en x=-1 la singularidad es removible
        \addplot[domain=-1.8:2.6, samples=400] {(x^5-2*x^2+3)/(x^3+1.00001)};
        \node[font=\small] at (axis cs:0.75,4.4) {$f(x)=\dfrac{x^{5}-2x^{2}+3}{x^{3}+1}$};
        \node[font=\small] at (axis cs:-1.3,1.1) {$y=x^{2}$};
    \end{axis}
\end{tikzpicture}
\end{document}
